\section{De los LLM a lo multimodal: la visión de cómputo unificado de NVIDIA}

\subsection{La hipótesis central de Jensen: la escalabilidad aún no llega a su límite}

En la entrevista, Jensen insiste en «si ya puede hacer esto, ¿hasta dónde más puede llegar?». Detrás de esto está la hipótesis de que:
\begin{itemize}
    \item Seguir aumentando la escala (modelo, datos, cómputo) sigue teniendo retorno.
    \item La arquitectura irá iterando, pero la necesidad de «paralelismo de alto rendimiento + interconexión de alto ancho de banda + coordinación de la pila de software» no desaparecerá.
\end{itemize}

Esta hipótesis explica por qué NVIDIA apuesta a la vez por chips, interconexión, software de sistemas y cadena de herramientas de modelos, en lugar de enfocarse en un solo componente.

\begin{importantbox}{Una visión unificada de la tarea: en esencia, mapeo de tokens}
En la entrevista, Jensen menciona el paso de texto a imagen, de texto a tokens de acción, y de secuencias biológicas a predicción de estructuras. Su visión unificada es que las distintas modalidades pueden convertirse en una representación tokenizada y luego pasar por el mismo tipo de sistema de aprendizaje para realizar el mapeo.

Esto no significa que las diferencias entre tareas desaparezcan, sino que muestra que el «núcleo de cómputo reutilizable» se está ampliando.
\end{importantbox}

\subsection{Caso: la IA retroalimenta el renderizado gráfico}

En la parte final, Jensen presenta un caso clave: en una escena en 4K no se calcula cada píxel de manera tradicional, sino que se calcula una parte de los píxeles con alta calidad y se deja que la IA prediga y complete el resto, logrando una optimización conjunta de calidad y rendimiento.

\begin{figure}[H]
\centering
\includegraphics[width=0.88\textwidth]{frame-geforce-b.jpg}
\caption{La idea de renderizado que calcula una parte de los píxeles y usa IA para completar el resto \protect\footnotemark}
\end{figure}
\footnotetext{Intervalo de tiempo del video: 00:53:14--00:53:27.}

\begin{knowledgebox}{La estructura de interés compuesto de «la IA retroalimenta a la IA»}
GeForce alguna vez ayudó a la comunidad de investigación a entrenar AlexNet; y ahora la IA regresa al renderizado gráfico y a la optimización de sistemas. Este tipo de ciclo tecnológico genera un efecto de interés compuesto:
\begin{itemize}
    \item La aplicación de una generación amplía la escala del hardware.
    \item Una escala de hardware mayor, a su vez, sostiene a la siguiente generación de modelos.
    \item Los nuevos modelos vuelven a retroalimentar y optimizar las aplicaciones originales.
\end{itemize}
\end{knowledgebox}

\subsection{Resumen de la sección}
La narrativa de mediano y largo plazo de NVIDIA no es que «gane un modelo grande en particular», sino que, a medida que los sistemas de aprendizaje multimodal sigan expandiéndose, el margen de reutilización de la plataforma de paralelismo general seguirá ampliándose.

